\chapter{Supplement G, Part 1: The Gaussian Realization and Its Geometry}
\label{comp:sg1}

Supplement G asks what the single-point theory becomes when every belief and model law is
multivariate Gaussian (\SG{ch:gaussian-introduction}). This chapter covers its introduction and its
first five chapters; aggregation, Kron reduction, nonflat links, the operator arrow, and reciprocal
pairs are in \cref{comp:sg2}. Three facts frame what follows. The supplement is a realization, not the
ambient theory: it asserts nothing for a non-Gaussian belief fiber beyond the exponential-family
statements of G1, and the main manuscript uses none of its results. The Gaussian lives on real fibers,
\(\R^K\) for beliefs and \(\R^{d_m}\) for model presentations with \(K\neq d_m\) allowed, and the
structure group reaches them only through its represented image \(\Gamma_b=\rho_b(G)\subseteq\GL(K)\);
some statements hold for every \(\Gamma_b\), and a named set is proved under the specialization
\(\Gamma_b=\GL^+(K)\) (\SG{hyp:gen-lg-frame-group}). And the belief precision analyzed throughout is
that of the fast belief at frozen model latents, \(p(k\given m,o,X)\), not the joint belief--model
precision.

Most of the mathematics is finite-dimensional linear algebra and Gaussian integration, much of it
textbook in substance. What is specific to the program is the typing (which objects are laws, which are
operators, which are artifacts of a chart or frame) and the analysis of when a directed linear-Gaussian
network yields the precision a physicist would write first, a matrix-weighted graph Laplacian plus local
mass terms. Only \SG{prop:exp-cumulants-fisher} is tagged STANDARD; every other statement in G1 to G5 is
ESTABLISHED with a proof, and none rests on computation. In one line each: G1 shows that boundary
points and operator fixed points carry no law; G2 gives the information form, the forest precision, and
one Schur lemma behind every later marginal; G3 shows that the Laplacian-plus-self-term shape is a
declared restriction and says when it holds, exactly in one chart and at law level on forests and pairs;
G4 prices recognition restrictions and finds that
none selects a partition; G5 separates law-level invariants from chart and frame artifacts.

\section{Exponential families and the operator layer (G1)}
\label{comp:sg1-expfamily}

\subsection{The question and the answer}

Later chapters apply algebraic maps (aggregation, rescaling, Schur complementation) to matrices in a
cone, and a matrix in the cone need not be the precision of any normalizable law. \SG{ch:expfamily} asks
what happens when a map of natural parameters leaves the natural domain \(\mathsf N\). On its boundary
nothing is left: log partition, mean parameter, Fisher metric, relative entropies, and evidence bound
diverge or are undefined (\SG{prop:exp-domain-boundary-no-law}); a map can keep every iterate inside
while its only fixed point is on the boundary (\SG{prop:exp-fixed-point-no-law}); and pinning,
quotienting, and keeping a mass coordinate are three different repairs (\SG{sec:exp-repairs}). Inside,
relative entropy is a Bregman divergence (\SG{prop:exp-kl-bregman}) and affine mean constraints have a
unique information projection (\SG{prop:exp-affine-mean-projection}).

\subsection{The picture}

\begin{analogy}{exact}
An exponential family is a Gibbs ensemble: the statistic \(\mathsf T\) lists observables, the natural
parameter \(\lambda\) their conjugate fields, \(\mathsf A(\lambda)\) is a Massieu function,
\(\nabla\mathsf A\) the thermodynamic averages, and \(\nabla^2\mathsf A=\operatorname{Cov}(\mathsf T)\)
the fluctuation--response relation, which is also the Fisher metric (\SG{prop:exp-cumulants-fisher}). For
the Gaussian, \(\mathsf T=(y,yy^\top)\) and \(\lambda=(\mathsf h,-\tfrac12\Lambda)\), a source and a
coupling matrix. The information projection is the family-restricted form of minimum relative entropy
\citep{jaynes2003probability}.
\end{analogy}

The boundary results are an infrared statement. For the scalar Gaussian with \(\mathsf T(z)=-z^2/2\),
\(\lambda\) is the stiffness and \(\mathsf N=(0,\infty)\); the map \(\lambda\mapsto\lambda/2\) keeps
every iterate normalizable, but its limit \(\lambda=0\) is the flat weight on the line
(\SGeq{eq:exp-scale-witness}), inside the operator cone \([0,\infty)\) and outside the probability
layer.

\begin{analogy}{structural}
This is the massless limit of a free field, whose zero mode has no normalizable measure. Physicists fix a
boundary value, remove the zero mode, or keep a small mass as a running coupling; G1's pinning
(\SG{def:exp-pinning}), quotienting (\SG{def:exp-quotienting}), and mass sector
(\SG{def:exp-mass-sector}) are these three repairs, and they give three different objects. A
pseudoinverse and pseudodeterminant in the full-space density formula compute the quotient normalizer,
on a smaller space with another reference measure, while the full-space integral stays infinite
(\SG{prop:exp-quotient-normalization}).
\end{analogy}

\subsection{Objects and notation}

\begin{decoder}
Canonical family \((\nu,\mathsf T,\lambda,\mathsf A,\mathsf N)\) & Laws \(\propto e^{\langle\lambda,\mathsf T\rangle}\) relative to \(\nu\); \(\mathsf N\) is where \(\mathsf A<\infty\) & Gibbs ensemble\\
Regular, minimal & \(\mathsf N\) open; no affine redundancy, so parameters determine laws & No redundant coordinates\\
Operator cone \(\mathcal W\); probability layer \(\mathcal W\cap\mathsf N\) & Parameters preserved by the maps; the part that denotes laws & Couplings with unphysical points; normalizable couplings\\
Bregman divergence; information projection & Error of the linear extrapolation of \(\mathsf A\); closest member under affine mean constraints & Availability; maximum entropy\\
\end{decoder}

Traps: the sans-serif \(\mathsf A\) is the log partition, the italic \(A_i\) of later chapters are self
terms; the bold \(\boldsymbol\mu\) is a mean parameter, for the Gaussian the pair \((\mu,M_2)\) of mean
and second moment (\SG{def:gauss-coordinate-charts}); ``operator'' means a matrix acted on by an
algebraic map, nothing quantum.

\subsection{The formulas, explained}

Inside the domain,
\[
\KL(Q_{\lambda'}\Vert Q_\lambda)=\mathsf A(\lambda)-\mathsf A(\lambda')-\langle\nabla\mathsf A(\lambda'),\lambda-\lambda'\rangle
\]
(\SG{prop:exp-kl-bregman}): the difference of log partitions minus the first-order change predicted from
\(\lambda'\). Convexity makes it nonnegative; for nearby parameters it is the Fisher quadratic, and at
finite separation it is not, a difference the supplement insists on. For a closed affine set of mean
parameters meeting the interior, the projection \(Q_{\widehat\lambda}\) is unique and every
constraint-satisfying member obeys
\[
\KL(Q_{\lambda'}\Vert Q_\lambda)=\KL(Q_{\lambda'}\Vert Q_{\widehat\lambda})+\KL(Q_{\widehat\lambda}\Vert Q_\lambda),
\]
as a squared distance splits under orthogonal projection (\SG{prop:exp-affine-mean-projection}); the
Bernoulli examples of \SG{rem:exp-projection-hypotheses} show that interior feasibility, closedness,
and convexity are each needed. The scope paragraph is a real warning: no coarse-graining statement
follows from exponential-family structure alone, since an operator arrow needs, among other things, a
declared rescaling, and a law-level arrow needs finite normalizers at every scale.

\section{The Gaussian realization at the base point (G2)}
\label{comp:sg1-realization}

\subsection{The question and the answer}

\SG{ch:generative-gaussian} asks what the directed law becomes when every kernel is linear Gaussian, and
what the precision blocks mean. The precision is at once the coupling matrix and the Fisher metric of the
mean (\SG{prop:gauss-log-normalizer-moments}); a zero block is conditional independence, not zero
covariance (\SG{prop:gauss-conditional-precision}); the marginal precision of a set of directions is a
Schur complement below the restricted precision, with equality exactly under \(\Lambda\)-orthogonality
(\SG{lem:gauss-schur}). A forest assembles into an explicit sparse precision
(\SG{prop:gauss-assembled-precision-blocks}), every kernel is frame covariant
(\SG{prop:gen-lg-kernel-covariance}), and the learning step has a closed form that fails to exist in two
named cases (\SG{prop:gen-lg-m-step}).

\subsection{The picture}

\begin{analogy}{exact}
In \(q\propto\exp(\mathsf h^\top y-\tfrac12y^\top\Lambda y)\), \(\Lambda\) is the quadratic part of a
lattice action (the inverse propagator), \(\mathsf h\) an external source, and \(\Sigma=\Lambda^{-1}\)
the propagator. A zero off-diagonal block means no direct bond while the propagator still connects the
sites through paths (\SGeq{eq:gauss-precision-not-covariance}). The Schur complement is integrating out,
exact at tree level for a Gaussian. The inequality \(\Lambda_{bb}^{-1}\preceq(\Lambda^{-1})_{bb}\)
(\SG{prop:restrict-optimum-vs-marginal}) is the Le Chatelier--Braun inequality: a block fluctuates less
with the rest clamped than with the rest free to relax, which is why a factorized recognition law, whose
block covariances are the clamped ones, underestimates every marginal fluctuation.
\end{analogy}

The directed model (\SG{def:gen-lg-forest}) is a state-space model on a tree: each nonroot belief is its
parent's belief carried by the link \(\Theta_{ij}\), multiplied by a gain \(F_i\), shifted by a bridge
term \(\mathsf B_im_i\) and an offset, plus noise of precision \(\Lambda_i\), with linear-Gaussian
observations. The product \(\Omega_i=F_i\Theta_{i\operatorname{par}(i)}\) is the transition. At frozen
model latents the belief sector is the forest
\(k_i\given k_{\operatorname{par}(i)}\sim\mathcal N(\Omega_ik_{\operatorname{par}(i)},\Lambda_i^{-1})\),
and model latents, offsets, and the model part of observations move only \(\mathsf h\): model
presentations act on beliefs as an external source, never as a coupling. This is the Born--Oppenheimer
move of the main manuscript's fast--slow split (\MT{sec:obs-fast-slow}), cleaner than usual because the
fast precision does not depend on the frozen variables at all.

\subsection{Objects and notation}

\begin{decoder}
\((\mathsf h,\Lambda)\); \(\mu=\Lambda^{-1}\mathsf h\), \(\Sigma=\Lambda^{-1}\) & Information vector, precision; mean, covariance & Source, couplings; propagator\\
Four charts & Information \((\mathsf h,\Lambda)\), natural \((\mathsf h,-\tfrac12\Lambda)\), expectation \((\mu,M_2)\), moment \((\mu,\Sigma)\) & Choices of thermodynamic variables\\
\(F_i\), \(\Theta_{ij}\), \(\Omega_i\) & Gain, represented link, transition & Transfer matrix times transporter\\
Restricted versus marginal precision & \(E^\top\Lambda E\) versus \((E^\top\Lambda^{-1}E)^{-1}\) & Clamped versus relaxed stiffness\\
M-coordinate & Maximizing expected complete log likelihood over generative parameters & Learning (EM's M-step)\\
\end{decoder}

Traps: one subscript on \(\Lambda\) is a kernel's innovation precision, two a block of the assembled
matrix (the \(W_i\) of G3 follows the same rule); symmetric-matrix coordinates scale off-diagonal entries
by \(\sqrt2\); the moment chart is not Legendre dual to the natural chart, the expectation chart is;
``pullback'' means restriction along an injective parameter map.

\subsection{The formulas, explained}

The log normalizer
\(\mathsf A(\mathsf h,\Lambda)=\tfrac12\mathsf h^\top\Lambda^{-1}\mathsf h-\tfrac12\log\det\Lambda+\tfrac n2\log2\pi\)
is a source energy mediated by the propagator, a fluctuation determinant, and a measure constant; the
Fisher information in \(\mu\) at fixed \(\Sigma\) is \(\Lambda\) (\SG{prop:gauss-log-normalizer-moments}),
so the stiffness of the energy along a direction is the distinguishability of mean shifts along it. The
conditional mean \(\E[Y_b\given Y_{-b}]=\mu_b-\Lambda_{bb}^{-1}\sum_{c\neq b}\Lambda_{bc}(y_c-\mu_c)\),
with conditional covariance \(\Lambda_{bb}^{-1}\), is the local equilibrium of block \(b\) in its
neighbors' field, and as an update the coordinate-ascent step of G4. With an agreement coupling
\(\Lambda_{bc}=-W_{bc}\preceq0\), a neighbor's deviation pulls the block toward it in the \(\Lambda_{bb}\)
inner product but not necessarily in the Euclidean one (\SG{prop:gauss-conditional-precision}, part
(iii)): agreement is relative to the
receiver's own metric, the only one available under general linear frames.

The forest precision has the blocks
\[
\Lambda_{i\operatorname{par}(i)}=-\Lambda_i\Omega_i,
\qquad
\Lambda_{ii}=\Lambda_i+\sum_{i'\in\operatorname{ch}(i)}\Omega_{i'}^\top\Lambda_{i'}\Omega_{i'},
\]
and zero between nonadjacent agents (\SG{prop:gauss-assembled-precision-blocks}): each factor puts its
noise stiffness on the child, a bond \(-\Lambda_i\Omega_i\) to the parent, and a back-reaction on the
parent. The bond need not be symmetric, which is the seed of G3. The belief--model cross block is
\(K\times d_m\), so no family of square agent blocks hosts the joint precision when \(K\neq d_m\).
Frame covariance (\SG{prop:gen-lg-kernel-covariance}) is lattice-gauge covariance of a matter action,
with gains transforming by conjugation, covariances by congruence, and observation maps absorbing the
inverse frame; it holds for every \(\Gamma_b\), and for \(\Gamma_b\subseteq\operatorname{O}(K)\) ordinary
eigenvalues are also invariant (\SG{prop:gen-lg-compact-frames}), which fails under \(\GL^+(K)\).

The learning step (\SG{prop:gen-lg-m-step}) is least squares plus residual covariance,
\(\E\log\mathcal N(t;Fw,\Sigma)\leq-\tfrac12\log\det(2\pi e\Sigma_*)\), attained when
\(\Sigma_*\succ0\); an exact regression in some direction makes the likelihood unbounded, the
variance-collapse degeneracy of maximum likelihood. The named cases are a record factor at one observed
value and a link learned inside the invertible group, whose supremum can sit at a singular map outside
the group: for \(K=1\) with independent centered \(k_i,k_j\), the optimized term
\(-\tfrac12\log(2\pi e(\E k_i^2+\mathsf t^2\E k_j^2))\) increases as \(\mathsf t\to0\) and is never
attained. Noncompactness of the link group can make learning ill posed; the scalar witness says nothing
about frequency at large \(K\). The communication formulas (\SG{prop:geo-selective-lift},
\SG{prop:obs-selective-gaussian-record}) are Gaussian conditioning and the Gaussian channel; the record's
mutual information is gauge invariant and at rank one takes every value in \((0,\infty)\).

\section{The interaction family (G3)}
\label{comp:sg1-interaction}

\subsection{The question and the answer}

The coarse-graining results of the second half act on one shape of precision, the interaction family
\(\mathcal I(V)\) (\SG{def:gauss-interaction-family}):
\[
\Lambda_{ii}=A_i+\sum_{j\neq i}W_{ij},
\qquad
\Lambda_{ij}=-W_{ij}\quad(i\neq j),
\qquad
A_i\succeq0,\quad W_{ij}=W_{ji}\succeq0.
\]
\SG{ch:gaussian} asks when the belief precision of G2 has this shape. It is a restriction, not a
consequence of linear-Gaussian directedness (\SG{prop:gauss-interaction-family-thin},
\SG{prop:gauss-global-condition-witnesses}). With
orientation-preserving frame changes, positive-determinant transitions on a forest can be gauged away and
a pair can be brought into the family exactly when \((-1)^K\det\Lambda_{12}>0\)
(\SG{prop:gauss-gain-absorption}). A reversing gain stays outside in every chart reached by
orientation-preserving frames, and a cycle with the wrong determinant sign, or for odd \(K\) the
three-agent v-structure, in every chart reached by any frames (\SG{prop:gauss-law-level-witnesses}). In the
link-trivializing chart two computable conditions decide membership
(\SG{prop:gauss-global-interaction-condition}), and there the family is adopted as a standing hypothesis
(\SG{hyp:gauss-global-interaction}).

\subsection{The picture}

A member's quadratic form is
\[
y^\top\Lambda y=\sum_{i}y_i^\top A_iy_i+\sum_{i<j}(y_i-y_j)^\top W_{ij}(y_i-y_j)
\]
(\SG{prop:gauss-interaction-energy-kernel}): \(K\)-component site variables, matrix springs
\(W_{ij}\) penalizing disagreement, and pinning or mass terms \(A_i\). Consensus configurations cost only
\(v^\top(\sum_iA_i)v\); they are the zero modes of the Laplacian part \(L\), lifted by self terms alone,
and at \(A=0\) they are the non-normalizable zero modes of G1. A singular bond leaves disagreement free
along its kernel, so \(\ker L\) can exceed the consensus subspace on a connected graph. With definite
bonds under flat comparison, one definite self term per connected component makes \(\Lambda\succ0\)
(\SG{prop:gauss-interaction-nonempty}).

\begin{analogy}{exact}
For \(K=1\), \(W_{ij}=J\) on lattice bonds and \(A_i=m^2\), a member is the massive lattice Laplacian
\(-\Delta+m^2\) of the lattice Gaussian free field; for general \(K\) it is the matrix-weighted Laplacian
of the consensus literature \citep{Trinh2018MatrixWeighted}. With nontrivial links the agreement energy
\(\sum(k_i-\Theta_{ij}k_j)^\top W(k_i-\Theta_{ij}k_j)\) is the covariant lattice Laplacian of a matter field
in a background gauge field \citep{Creutz1983}: the connection Laplacian of \citet{SingerWu2012} for
orthogonal links, a sheaf Laplacian \citep{HansenGhrist2019} for general linear ones. The program's links
are independent edge variables, as lattice links are, so this is exact at the level of the quadratic
form. G3's family is what the covariant Laplacian becomes in the pure-gauge sector: when every loop
product is the identity the links are a coboundary \(\Theta_{ij}=U_iU_j^{-1}\)
(\MT{prop:gauss-coboundary-trivialization}), and \(z_i=U_i^{-1}k_i\) turns every residual into
\(U_i(z_i-z_j)\) (\SG{hyp:gauss-flat-comparison}), giving an ordinary matrix-weighted Laplacian with
weights \(U_i^\top WU_i\). With a nontrivial Wilson loop no gauge does this; that case is excluded here
and treated in \cref{comp:sg2}.
\end{analogy}

\subsection{Objects and notation}

\begin{decoder}
\(\mathcal I(V)\), \(A_i\), \(W_{ij}\), \(L\) & Interaction family; self term; edge weight; Laplacian part & Massive matrix-weighted Laplacian; mass; bond\\
Link-trivializing chart & Coordinates \(z_i=U_i^{-1}k_i\) under flat comparison, unique up to one right factor per component & Gauge setting every link to one\\
Block reframing \(\mathbf P=\operatorname{diag}(\mathbf P_i)\) & Frame change agent by agent, \(\Lambda\mapsto\mathbf P^\top\Lambda\mathbf P\) & Site-dependent passive gauge transformation\\
Gain absorption; orientation criterion & Frames turning every transition into the identity; \((-1)^K\det\Lambda_{12}>0\) & Maximal-tree gauge; ferromagnetic bond sign (\(K=1\))\\
Chart-independent witness; v-structure & Outside the family in every admissible chart; a child with two parents & Frustrated plaquette; explaining away\\
\end{decoder}

Traps: \(\Omega_i\) is a transition, not a holonomy or frequency; \(W_i:=\Lambda_i\Omega_i\) with one
subscript is the candidate weight of the edge to the parent, \(W_{ij}\) an edge weight; ``flat'' means
trivial loop products of links at the one base point, with no base manifold in sight; an ``admissible''
chart is one reached by block reframings in the group a statement names.

\subsection{The formulas, explained}

Why the shape is not automatic is visible in one block. A receiving factor contributes the bond
\(-\Lambda_i\Omega_i\), and membership needs it to be \(-W\) with \(W\) symmetric positive
semidefinite, which nothing in a linear-Gaussian model forces. Even with links removed the gain remains:
\(\widetilde\Lambda_i=\begin{psmallmatrix}2&1\\1&3\end{psmallmatrix}\) with the quarter-turn
\(\widetilde\Omega_i=\begin{psmallmatrix}0&1\\-1&0\end{psmallmatrix}\) gives the bond
\(\begin{psmallmatrix}1&-2\\3&-1\end{psmallmatrix}\), and \(F_i=-I\) gives a symmetric bond of the wrong
sign (\SG{prop:gauss-global-condition-witnesses}). In a fixed chart every off-diagonal block must be
symmetric, \(\binom N2K(K-1)/2\) linear conditions, so the family is Lebesgue-null once \(K,N\geq2\)
(\SG{prop:gauss-interaction-family-thin}).

Thinness concerns one chart; the law-level question is whether some frame change puts the precision into
the family, and the answer is gauge-theoretic. A reframing \(k_i=\mathbf P_iz_i\) sends \(\Omega_i\) to
\(\mathbf P_i^{-1}\Omega_i\mathbf P_{\operatorname{par}(i)}\), the transformation law of a link
(\SG{prop:gauss-gain-absorption}, part (o)): gain times link behaves as one link, and no gauge
transformation can tell a gain from a link. On a tree every link configuration is pure gauge, so frames
\(\mathbf P_r=I\) at roots and \(\mathbf P_i=\Omega_i\mathbf P_{\operatorname{par}(i)}\) turn every
residual into \(z_i-z_{\operatorname{par}(i)}\), leaving pure agreement bonds
\(W=\mathbf P_i^\top\Lambda_i\mathbf P_i\) and self terms only at roots: the maximal-tree gauge applied to
transitions. The frames must lie in the frame group, which under \(\GL^+(K)\) requires
\(\det\Omega_i>0\); for a forest with invertible transitions this is also necessary (part (iii)). For a
pair the only obstruction is the sign
\[
(-1)^K\det\Lambda_{12}>0
\]
(part (ii)). Informally: a member with invertible bond has \(\Lambda_{12}=-W\), \(W\succ0\), so
\((-1)^K\det\Lambda_{12}=\det W>0\), and a reframing multiplies \(\det\Lambda_{12}\) by
\(\det\mathbf P_1\det\mathbf P_2>0\). The sign is a \(\mathbb Z_2\) gauge invariant of the edge. Admitting
orientation-reversing frames, \(\Gamma_b=\GL(K)\), removes it (part (iv)).

\begin{analogy}{exact}
For \(K=1\) this is the gauge structure of Ising spin glasses \citep{mezard1987spin}. The family is the
set of precisions with ferromagnetic couplings \(\Lambda_{ij}\leq0\) and nonnegative self terms; positive
rescalings cannot flip a bond sign (the orientation criterion); \(k_i\mapsto-k_i\) is the Mattis gauge
transformation, which makes any sign pattern on a tree ferromagnetic; and on a loop the product of bond
signs is invariant (frustration). \SG{prop:gauss-law-level-witnesses}(ii) is the \(K\times K\) version:
the sign of the product of block determinants around a cycle of \(n_\circ\) agents is invariant under
every \(\GL(K)\) reframing, and membership forces it to be \((-1)^{Kn_\circ}\). Part (iii) reads explaining
away as frustration: two parents of one child are coupled through
\(\Lambda_{jj'}=\Omega_j^\top\Lambda_i\Omega_{j'}\), and the triangle has the wrong determinant sign for
every odd \(K\). The correspondence is exact at the level of these signs and says nothing beyond them.
\end{analogy}

In the link-trivializing chart membership is decided exactly
(\SG{prop:gauss-global-interaction-condition}): each \(W_i\) must be symmetric positive semidefinite, and
each self term
\[
A_i=\Lambda_i-W_i+\sum_{i'\in\operatorname{ch}(i)}\left(W_{i'}\Lambda_{i'}^{-1}W_{i'}-W_{i'}\right)
\]
must be positive semidefinite (no \(-W_i\) at a root), with \(A_i+O_i\succeq0\) once observation
precisions are added. \(\Lambda_i-W_i\) is what a factor leaves on the child after paying for its bond,
and \(W_{i'}\Lambda_{i'}^{-1}W_{i'}-W_{i'}\) what each child's factor leaves on the parent. For a scalar
transition \(\omega\) these are \(\Lambda(1-\omega)\) and \(\Lambda\omega(\omega-1)\): a contracting gain
leaves a negative self term on the parent, an expanding one on the child, and only \(\omega=1\) is safe
factor by factor (\SG{prop:gauss-edge-local-characterization}, \SG{cor:gauss-invertible-gain-factor}).
Globally a root prior or observation can pay for a negative contribution, so factorwise membership is
strictly stronger than the hypothesis (\SG{hyp:gauss-factorwise-agreement}). A physicist can read the
condition as a ban on net tachyonic self terms.

Two consequences need emphasis. First, the adopted hypothesis is stronger than law-level membership. It
requires membership in the link-trivializing chart, and the scope paragraph says nothing selects another
chart. The prior of part (iii) of \SG{prop:gauss-global-interaction-condition}, with transition
\(\tfrac12\), is outside the family in that chart and, by gain absorption, inside it in its gain-absorbing
chart. The choice matters downstream: aggregation identifies agents by equal coordinates in the chosen
chart, so the chart decides whether agreement means link-transported equality \(k_i=\Theta_{ij}k_j\) or
equality with the prediction \(k_i=\Omega_ik_{\operatorname{par}(i)}\). That is a modeling decision, and
the supplement makes it explicitly. Second, membership is not invariant under general block reframings,
which preserve inertia but not the symmetry of off-diagonal blocks
(\SG{rem:gauss-membership-frame-dependence}); frame-independent statements need the invariants of G5.

\section{Recognition restrictions (G4)}
\label{comp:sg1-restrictions}

\subsection{The question and the answer}

Fix a Gaussian model with posterior \(p^\star=\mathcal N(\mu,\Lambda^{-1})\), \(\Lambda\succ0\). By the
restriction principle (\MT{prop:restrict-principle}), shrinking the recognition family lowers the
optimized bound by the rise in minimal relative entropy to the posterior. \SG{ch:restrictions} computes
that loss for covariance factorization (the determinant gap, \SG{thm:restrict-determinant-gap}), for mean
restrictions (\SG{prop:restrict-mean-cost}), and for identifying agents (\(+\infty\),
\SG{prop:cg-epsilon-identification-cost}), and shows that the gap selects no structure
(\SG{prop:restrict-gap-no-selection}).

\subsection{The picture}

For a Gaussian candidate,
\[
\KL(Q\Vert p^\star)=\tfrac12\left[\Tr(\Lambda\Sigma_Q)+(\mu_Q-\mu)^\top\Lambda(\mu_Q-\mu)-n-\log\det\Lambda-\log\det\Sigma_Q\right]
\]
(\SG{prop:restrict-gaussian-reverse-kl}). The mean enters only the quadratic term and the covariance only
the trace and log determinant, so mean and covariance restrictions are separate problems whose costs add
(\SG{cor:restrict-combined-costs-add}).

\begin{analogy}{exact}
The bound's gap is the slack in the Gibbs--Bogoliubov--Feynman inequality. For a Gaussian field with a
block-factorized trial law, the optimal means are exact and the optimal block covariances are the inverse
on-block stiffnesses \(\Lambda_{bb}^{-1}\) (\SG{thm:restrict-exact-block-optimum}): Weiss mean-field
theory, which drops the feedback of fluctuations through the environment, with the gap as its
free-energy error \citep{Bishop2006,Wainwright2008}. The mean-restriction cost
\(\tfrac12\mu_\perp^\top(B_\perp^\top\Lambda^{-1}B_\perp)^{-1}\mu_\perp\) is the free-energy cost of
constraining a Gaussian field's mean: half the squared displacement weighted by the inverse relaxed
susceptibility. The bare stiffness \(B_\perp^\top\Lambda B_\perp\) overstates it unless the constrained
directions are \(\Lambda\)-orthogonal to the rest (\SGeq{eq:cg-loewner}).
\end{analogy}

For two blocks the gap is their mutual information, \(-\tfrac12\sum_a\log(1-\check\rho_a^2)\) in the
canonical correlations (\SG{prop:restrict-canonical-correlation}); for more blocks it is the dual total
correlation \(h(Y)-\sum_bh(Y_b\given Y_{-b})\) of \citet{Han1978}
(\SG{cor:restrict-gap-dual-total-correlation}), built from the conditional entropies that
\(\Lambda_{bb}\) encodes, not the total correlation built from marginal ones (for three unit scalars with
pairwise correlation \(\tfrac12\), \(\tfrac12\log\tfrac{27}{16}\) against \(\tfrac12\log2\)). The gap is
observation free because a data-independent joint precision makes conditioning on \(o\) move only the
posterior mean (\SG{cor:restrict-gap-observation-free}); the mean cost is quadratic in the data
(\SG{cor:restrict-mean-cost-data}).

Identification of agents is the case that matters for coarse-graining. Let \(\operatorname{range}S\) be
the cluster-constant configurations of a hard aggregation, of dimension \(r=n_cK<n\). A recognition law
making all agents of a cluster literally equal lives on that subspace, to which the posterior, having a
Lebesgue density, gives probability zero. Absolute continuity fails (clause (H3) of
\MT{hyp:elbo-evidence-domain}), and the relative entropy is \(+\infty\). Thickening by a variance
\(\varepsilon\) in each identified direction gives
\[
\mathbb D_\varepsilon=\mathbb D_{\rm tie}+\frac{n-r}2\log\frac1\varepsilon-\frac{n-r}2
+\frac12\log\det(B_\perp^\top\Lambda^{-1}B_\perp)+\frac\varepsilon2\Tr(B_\perp^\top\Lambda B_\perp)
\]
(\SG{prop:cg-epsilon-identification-cost}), one \(\tfrac12\log(1/\varepsilon)\) per frozen continuous
coordinate, as a delta-function constraint costs entropy logarithmic in the resolution. Identification is
therefore a generative construction, the fine law conditioned on zero within-cluster disagreement
(\SG{prop:cg-aggregation-conditioning}); at finite cost a recognition law can only tie its means, the
mean-tie family with cost \(\mathbb D_{\rm tie}\) (\SG{prop:cg-mean-tie-cost}). Restricting the energy to
\(\operatorname{range}S\) and renormalizing, the ``hard trace,'' changes the evidence
(\SG{prop:restrict-hard-trace-evidence}); with a data-independent joint precision its posterior is the
image of the fine one under \(P_\Lambda=(S^\top\Lambda S)^{-1}S^\top\Lambda\), an evidence-preserving
coarse channel (\MT{thm:cg-evidence-preserving-channel}).

\subsection{Objects and notation}

\begin{decoder}
Recognition restriction; partition \(\mathcal P\) & Shrinking the candidate family; blocks with block-diagonal candidate covariance & Trial ensemble; mean-field clusters\\
Cost, gap & Loss of optimized bound, the rise of minimal relative entropy & Variational free-energy error\\
Mean tie; identification & Candidate means cluster constant; candidate supported on \(\operatorname{range}S\) & Constrained order parameter; delta-function constraint\\
Hard trace; \(P_\Lambda\) & Energy restricted to \(\operatorname{range}S\); the map carrying the fine posterior to it & Constrained partition function; Galerkin projection\\
\end{decoder}

Traps: ``restriction'' restricts the variational family, not a tensor; ``structure'' is the partition
used by the recognition law; the orthonormal basis \(B\) is unrelated to the bridge \(\mathsf B_i\);
``understates the covariance'' is said of the optimal factorized law, never of the model.

\subsection{The formulas, explained}

The determinant gap,
\[
\log p_\theta(o\given X)-\Lelbo^\star_{\mathcal P}=\frac12\left[\sum_{b\in\mathcal P}\log\det\Lambda_{bb}-\log\det\Lambda\right]\geq0,
\]
compares the product of on-block stiffness determinants with the full determinant; by Fischer's
inequality it vanishes exactly when every cross block does. When every block is a union of whole agents,
a block reframing shifts both terms by the same Jacobian, so the gap is invariant under every reframing,
whatever the frame group (\SG{prop:ig-frame-dependent-spectra-determinants}, part (iii)); a partition
that splits an agent is not.

Does any analyzed cost select a partition scale? No. The gap grows under refinement by an explicit
nonnegative amount (\SG{prop:restrict-refinement-monotonicity}) and vanishes at one block, so minimizing
it selects the unrestricted family (\SG{prop:restrict-gap-no-selection}). Non-nested partitions are
ordered by no structural feature: for
\(\Lambda=\begin{psmallmatrix}1&a&0\\a&1&b\\0&b&1\end{psmallmatrix}\) the partitions
\(\{\{1,2\},\{3\}\}\) and \(\{\{1\},\{2,3\}\}\) have equal block sizes and swap order when \(|a|\) and
\(|b|\) are exchanged (\SG{prop:restrict-nonnested-unordered}). The mean-tie cost vanishes at the finest
partition and identification is infinite for every nontrivial cluster. Each cost is extremized at an
endpoint, the Gaussian instance of \MT{open:restrict-intrinsic-criterion} and \MT{open:cg-natural-selector}.
One functional the chapter writes down is not examined in this role: the combined cost
\(\mathbb D_{\rm tie}+\mathbb D_{\rm fact}\) of \SG{cor:restrict-combined-costs-add} (mean-tie cost plus
cluster factorization gap, both displayed in \SG{prop:cg-mean-tie-cost}), with both
restrictions indexed by the same clusters, is the exact loss of one admissible family, and its terms move
in opposite directions under merging (the first cannot decrease because merging shrinks the allowed
means; the second cannot increase). The no-selection statements concern each cost separately; whether
the sum has a nondegenerate minimizer is not analyzed.

\section{Fisher geometry and gauge invariants (G5)}
\label{comp:sg1-fisher}

\subsection{The question and the answer}

\SG{ch:infogeometry} asks which geometric statements about the Gaussian family concern laws and which
concern a chart or frame. ``The Fisher mean block is the precision'' holds in the moment chart but not,
for coordinate directions, in the expectation chart (\SG{prop:ig-fisher-moment-chart},
\SG{prop:ig-fisher-expectation-chart}). Under frame changes, the generalized spectrum of the pencil
\((L,\Lambda)\) is invariant when regular (\SG{prop:grg-pencil-invariance}), while ordinary eigenvalues
and single subspace determinants are frame artifacts once the frame group contains the positive scalars
(\SG{prop:ig-frame-dependent-spectra-determinants}), and under \(\Gamma_b\supseteq\GL^+(K)\) a coupling
matrix keeps only its rank (\SG{prop:grg-residual-fiber-label}). Standard: the moment-chart metric of the
multivariate normal \citep{Skovgaard1984} with its affine isometries \citep{PineleStrapassonCosta2020},
exponential-chart duality \citep{AmariNagaoka2000}, congruence invariance and Sylvester's law
\citep{HornJohnson2013}. New: the typing of which quantities survive the represented frame group, and a
scoped statement of when a graph Laplacian is a Fisher tensor.

\subsection{The picture}

The two mean blocks differ because \(\partial/\partial\mu\) at fixed \(\Sigma\) and at fixed \(M_2\) are
different tangent vectors: the second must shrink the covariance to compensate
(\(\Sigma=M_2-\mu\mu^\top\)), and that extra motion has Fisher length, governed by the unbounded
Mahalanobis norm \(\mu^\top\Lambda\mu\). The quotient metric, a Schur complement, is exactly \(\Lambda\)
and sits below the restricted one, the clamped-versus-relaxed ordering of G2. As with \(C_p\neq C_V\), a
derivative depends on what is held fixed; an identity mixing mean and covariance sectors must name its
chart and say whether it restricts or quotients.

Under a general linear frame group there is no metric on the fiber, so the eigenvalues of one quadratic
form are chart data: rescaling all frames by \(\mathrm c>0\) multiplies every eigenvalue of \(\Lambda\) by
\(\mathrm c^2\), and a threshold on absolute eigenvalues can accept or reject any direction by choice of
\(\mathrm c\) (\SG{prop:ig-frame-dependent-spectra-determinants}, part (i)). Only relative quantities and
law functionals survive.

\begin{analogy}{exact}
Normal-mode frequencies are generalized eigenvalues of the potential form relative to the kinetic form
and do not depend on the generalized coordinates, whereas eigenvalues of the stiffness matrix alone do.
The pencil \(Lx=\mathsf d\Lambda x\) is this construction with the Laplacian part as potential and the
full precision as kinetic form; both transform by the same congruence, so the \(\mathsf d_k\) are
invariant. Each \(\mathsf d_k\in[0,1]\) is the fraction of precision along its direction that bonds supply
rather than self terms: consensus sits at \(0\), directions no self term touches at \(1\)
(\SG{prop:ig-generalized-spectrum-localization}).
\end{analogy}

\subsection{Objects and notation}

\begin{decoder}
Restriction versus quotient metric & Metric on coordinate directions versus orthogonal to a fiber & Clamped versus relaxed response\\
Block reframing \(\mathbf P\) & New-to-old frame change, \(\Lambda\mapsto\mathbf P^\top\Lambda\mathbf P\), \(\Sigma\mapsto\mathbf P^{-1}\Sigma\mathbf P^{-\top}\) & Site-dependent gauge transformation\\
Pencil \((L,\Lambda)\), regular & Pair of forms, \(Lx=\mathsf d\Lambda x\); determinant not identically zero & Normal-mode problem; no shared zero mode\\
Agent-aligned law functional; residual fiber label & Gap, mutual information, relative entropy of unions of agent blocks; congruence class of one semidefinite matrix & Gauge-invariant observable; inertia (Sylvester)\\
\end{decoder}

Traps: \(\mathsf d\) is a generalized eigenvalue, not a dimension; \(L\) is the Laplacian part, not the
bound \(\Lelbo\); a ``regular'' pencil and a ``regular'' exponential family are unrelated; the mass
\(\mathsf a\) below is a coefficient on a reference precision.

\subsection{The formulas, explained}

In the moment chart,
\[
g^F\bigl((\dot\mu_1,\dot\Sigma_1),(\dot\mu_2,\dot\Sigma_2)\bigr)=\dot\mu_1^\top\Lambda\dot\mu_2+\tfrac12\Tr(\Lambda\dot\Sigma_1\Lambda\dot\Sigma_2),
\]
and every affine map \((\mu,\Sigma)\mapsto(Q\mu+\mathrm v,Q\Sigma Q^\top)\) is an isometry; block
reframings act this way, so the Fisher geometry of the recognition family is gauge invariant for any
frame group. In the expectation chart
\(g_{[\mu\mu]}=(1+\mu^\top\Lambda\mu)\Lambda+\Lambda\mu\mu^\top\Lambda\), with eigenvalues relative to
\(\Lambda\) of \(1+\mu^\top\Lambda\mu\) (\(n-1\) times) and \(1+2\mu^\top\Lambda\mu\). The link between
Laplacian and Fisher geometry is narrow and exact:
\[
\KL\bigl(\mathcal N(\mu_i,W_{ij}^{-1})\Vert\mathcal N(\mu_j,W_{ij}^{-1})\bigr)=\tfrac12(\mu_i-\mu_j)^\top W_{ij}(\mu_i-\mu_j),
\]
summing over edges with definite weights to \(\tfrac12\mu^\top L\mu\) (\SG{prop:ig-fixed-covariance-kl}).
Outside Gaussian laws at that fixed covariance, or a frozen base law to second order, the chapter states
that a sample-space Laplacian is neither a recognition Fisher tensor nor a natural-gradient operator, a
scope limit on any use of a graph Laplacian as a preconditioner.

Pencil invariance is
\(\det(\mathbf P^\top L\mathbf P-\mathsf d\mathbf P^\top\Lambda\mathbf P)=\det(\mathbf P)^2\det(L-\mathsf d\Lambda)\).
The pencil is regular exactly when \(\ker L\cap\ker A=\{0\}\) (\SG{prop:grg-pencil-regularity}); at the
massless boundary \(A=0\) it is \((L,L)\), irregular, and says nothing about the graph until a pin,
quotient, or mass is declared. With a mass, \(\Lambda=L+\mathsf a\Lambda_{\mathrm{ref}}\) gives
\(\mathsf d=\mathsf d_{\mathrm{ref}}/(\mathsf d_{\mathrm{ref}}+\mathsf a)\)
(\SG{prop:grg-mass-pencil-transfer}), roughly the massive-propagator ratio \(p^2/(p^2+\mathsf a)\) on a
lattice; the small-eigenvalue counting exponent survives, and the limits \(\mathsf a\to0\) and
\(\mathsf d_{\mathrm{ref}}\to0\) do not commute (\SGeq{eq:grg-noncommuting-mass}). A subspace log
determinant picks up the Jacobian \(\log\det(B'^\top\mathbf P^\top\mathbf PB')\) (\SGeq{eq:ig-logdet-law}),
which orthogonal reframings kill and general linear ones do not; the gap of G4 cancels it for
agent-aligned blocks, while an agent-splitting gap moves from \(\tfrac12\log\tfrac43\) to
\(\tfrac12\log4\) in the \(N=1\), \(K=2\) witness. Under \(\Gamma_b\supseteq\GL^+(K)\) every semidefinite
matrix of rank \(r\) is congruent to \(\operatorname{diag}(I_r,0)\), so a common reframing leaves \(K+1\)
classes and no continuous modulus; under \(\operatorname{SO}(K)\) or \(\operatorname{O}(K)\) a definite
matrix keeps its \(K\) eigenvalues (\SG{prop:gen-lg-compact-frames}). With a noncompact frame group the
strength of a coupling matrix by itself is gauge; with a compact one it is physical.

One caution the wording does not make explicit. The pencil is formed after splitting \(\Lambda=L+A\) in
the chart where the family is imposed (\SG{def:ig-block-reframing}), the link-trivializing chart of
\SG{hyp:gauss-global-interaction}. Its spectrum is data of the law together with that split, hence with
its links, not of the law alone: a precision in the family in two inequivalent charts can have two
pencils. For \(N=2\), \(K=1\), \(\begin{psmallmatrix}2&-1\\-1&2\end{psmallmatrix}\) splits with \(W=1\),
\(A=(1,1)\) and has spectrum \(\{0,\tfrac23\}\), while its reframing by \(\operatorname{diag}(1,2)\),
\(\begin{psmallmatrix}2&-2\\-2&8\end{psmallmatrix}\), lies in the family with \(W=2\), \(A=(0,6)\) and
spectrum \(\{0,1\}\). This does not contradict \SG{prop:grg-pencil-invariance}, which concerns one
congruence of one pair; it limits \SG{rem:gauss-membership-frame-dependence}, which lists pencil
eigenvalues among frame-independent membership criteria, since forming the pencil presupposes a chart in
which membership already holds.

\section{What is proved, and how solid it is}
\label{comp:sg1-scorecard}

Every row is proved in the supplement or cited there to a checked source; there is no NUM statement and
no TOY positive result in this part. COND rows rest on a frame-group or graph hypothesis that fails in
cases the program cares about.

\begin{scorecard}
\multicolumn{4}{@{}l}{\textit{G1. Exponential families}}\\
\SG{prop:exp-domain-boundary-no-law} & No law, Fisher metric, or bound at a boundary point; an operator map can converge there (\SG{prop:exp-fixed-point-no-law}) & \SOLID & Negative; scalar witness suffices\\
\SG{sec:exp-repairs} & Pin, quotient, and mass sector are three distinct repairs & \CHOICE & Their two propositions are \SOLID\\
\SG{prop:exp-affine-mean-projection} & Cumulants and Fisher metric; relative entropy is Bregman; unique affine projection & \SOLID & Standard in substance; \SG{prop:exp-cumulants-fisher} is STANDARD\\
\multicolumn{4}{@{}l}{\textit{G2. Gaussian realization}}\\
\SG{hyp:gen-lg-frame-group} & Represented frame group, specialized to \(\GL^+(K)\) where named & \CHOICE & Not met by \(\operatorname{SU}(N)\) or a realified \(\GL(K,\mathbb C)\)\\
\SG{def:gen-lg-forest} & Directed linear-Gaussian forest & \CHOICE & General DAGs get one sentence\\
\SG{lem:gauss-schur} & Zero block is conditional independence; Schur complement below restricted precision & \SOLID & Standard in substance\\
\SG{prop:gen-lg-kernel-covariance} & Every kernel frame covariant for every \(\Gamma_b\); explicit forest blocks (\SG{prop:gauss-assembled-precision-blocks}) & \SOLID & Belief sector at frozen model latents\\
\SG{prop:gen-lg-m-step} & Closed-form learning step; no maximizer if the residual covariance is singular & \SOLID & Link witness is scalar: possibility, not frequency\\
\multicolumn{4}{@{}l}{\textit{G3. Interaction family}}\\
\SG{hyp:gauss-global-interaction} & The family, adopted in link-trivializing coordinates & \CHOICE & The chart fixes what agreement means\\
\SG{hyp:gauss-flat-comparison} & Belief links form a graph coboundary & \CHOICE & Automatic on forests; fails for generic links on cycles\\
\SG{prop:gauss-interaction-family-thin} & Null set in a fixed chart; not reframing invariant & \SOLID & Chart-level\\
\SG{prop:gauss-gain-absorption} & On forests, \(\GL^+(K)\) frames absorb positive-determinant transitions; pair criterion \((-1)^K\det\Lambda_{12}>0\); \(\GL(K)\) frames reach every invertible pair & \COND & Forest or pair; stated frame group; nothing sufficient for cycles or smaller groups\\
\SG{prop:gauss-law-level-witnesses} & Reversing gain, cycle determinant sign, odd-\(K\) v-structure exclude membership & \SOLID & Negative; (ii), (iii) for \(\GL(K)\); (iii) on three agents\\
\SG{prop:gauss-global-interaction-condition} & Exact membership test; prior balance sufficient, not necessary; factorwise stronger & \SOLID & Link-trivializing chart\\
\multicolumn{4}{@{}l}{\textit{G4. Recognition restrictions}}\\
\SG{thm:restrict-determinant-gap} & Exact mean, clamped covariances; loss is the gap, the dual total correlation (\SG{cor:restrict-gap-dual-total-correlation}) & \SOLID & Standard in substance\\
\SG{cor:restrict-gap-observation-free} & Gap independent of the observation; mean cost carries the data (\SG{cor:restrict-mean-cost-data}) & \SOLID & Needs data-independent joint precision\\
\SG{prop:cg-epsilon-identification-cost} & Identification costs \(+\infty\); regularized, \(\tfrac{n-r}2\log\tfrac1\varepsilon\) & \SOLID & Negative; absolute continuity\\
\SG{prop:restrict-gap-no-selection} & Gap monotone on chains, unordered off them, minimized by one block & \SOLID & Negative; combined cost not examined\\
\multicolumn{4}{@{}l}{\textit{G5. Fisher geometry and invariants}}\\
\SG{prop:ig-fisher-expectation-chart} & Mean block is \(\Lambda\) in the moment chart, larger in the expectation chart; quotient \(\Lambda\) & \SOLID & Standard in substance \citep{Skovgaard1984}\\
\SG{prop:ig-fixed-covariance-kl} & Edge relative entropy at covariance \(W_{ij}^{-1}\) is the Laplacian quadratic & \SOLID & Gaussian, fixed covariance only\\
\SG{prop:grg-pencil-invariance} & Regular pencil spectrum invariant; in \([0,1]\) & \SOLID & Invariant of law plus link-fixed split\\
\SG{prop:ig-frame-dependent-spectra-determinants} & Absolute spectra, subspace determinants frame dependent; agent gaps invariant & \SOLID & (i) needs positive scalars\\
\SG{prop:grg-residual-fiber-label} & A semidefinite coupling matrix is labeled by rank alone & \COND & Needs \(\Gamma_b\supseteq\GL^+(K)\); false for compact groups\\
\multicolumn{4}{@{}l}{\textit{Open}}\\
\SG{open:exp-mass-sector} & Mass-coordinate scaling; also \SG{open:exp-projection-commutation}, \SG{open:ig-model-space-fisher-transfer} & \OPEN & \Cref{comp:sg1-open}\\
\end{scorecard}

\begin{hygiene}
Conditions that can fail where the program lives: flat comparison (\SG{hyp:gauss-flat-comparison}),
without which the interaction hypothesis is not formulated; membership in the link-trivializing chart
(\SG{hyp:gauss-global-interaction}); the \(\GL^+(K)\) specialization (\SG{hyp:gen-lg-frame-group}),
needed by absorption and the rank-only label; frozen model latents; absolute continuity
(\MT{hyp:elbo-evidence-domain}, clause H3), which identification violates; a nondegenerate target and
regular pencil, which fail at the massless boundary \(A=0\); a data-independent joint precision; and a
nonsingular residual covariance in the learning step. Hygiene: standard Borel spaces, sigma-finite
references, measurability, regularity and minimality (automatic for the Gaussian), definite receiving
covariances (\SG{hyp:gen-lg-spd}), finite second moments, and orthonormal bases used as a computational
device.
\end{hygiene}

\section{Toy versus general, and what this part implies for \texorpdfstring{\(\GL(K)\)}{GL(K)} networks at large \texorpdfstring{\(N\)}{N}}
\label{comp:sg1-scale}

Nothing in G1 to G5 is TOY. Every theorem holds for arbitrary finite \(N\) and \(K\) (the two-agent
orientation criterion has a forest version at every \(N\)), and every small instance is a counterexample
to a universal statement, hence a solid negative result that shows a claim fails, not how often it fails
at \(N=512\). Thinness holds at every \(N\), and the cycle-sign condition
applies to any cycle with invertible blocks inside any network. Nothing here is a large-\(N\) statement
either: no thermodynamic limit, per-agent density, or scaling law.

\paragraph{The interaction family at scale.} Two conditions stand between a 512-agent Gaussian precision
with cycles and the family. Flat comparison requires every loop product of the belief links to be the
identity, automatic on a forest but violated by generic independent \(\GL(K)\) links on a graph with
many cycles; then the hypothesis behind aggregation and the operator arrow is not formulated, and the
nonflat case is in \cref{comp:sg2}. And the law-level criterion is complete only for forests and pairs;
for precisions whose graph has cycles, including those created by co-parents, only the necessary sign
condition is proved.

\paragraph{Frame groups.} For \(\GL(K,\R)\), \(\Gamma_b=\GL(K)\) contains \(\GL^+(K)\) and the positive
scalars, so every statement needing either holds. Part (iv) of \SG{prop:gauss-gain-absorption} reaches
every invertible pair, so the orientation obstruction disappears there, and the forest construction of
part (i) uses positive determinants only to keep its frames in \(\GL^+(K)\); the cycle-sign and
odd-\(K\) v-structure obstructions, stated for \(\GL(K)\) frames, remain. The gauge-invariant data of a Gaussian network under \(\GL(K,\R)\) are agent-aligned
law functionals, pencil spectra relative to the link-fixed split, cycle determinant signs, and holonomy
conjugacy classes (\MT{prop:geo-edge-twist-invariance}); absolute eigenvalues and single subspace
determinants are gauge artifacts. The supplement treats real Gaussians only. A complex group enters
through a real representation, for instance the realification on \(\R^{2K}\), whose image lies in
\(\GL^+(2K)\) and contains the positive scalars; statements proved for every \(\Gamma_b\) hold, the
rank-only label and the \(\GL^+\) sufficiency constructions do not apply as stated because their frames
need not be complex linear, and obstructions proved for larger groups still apply. For
\(\operatorname{SU}(N)\) and other compact groups, \(\Gamma_b\subseteq\operatorname{O}(K)\) after
choosing an invariant inner product: spectra are invariant, coupling matrices keep continuous eigenvalue
moduli, and no absorption result is stated.

\paragraph{Where compactness and orthogonality enter.} No theorem here assumes compactness. It is
excluded by the \(\GL^+(K)\) specialization, brings extra invariants when present, and its absence is
what the scalar link witness exploits (informally, on a compact group a continuous objective attains its
maximum, a standard fact the supplement does not state). A Euclidean structure enters as a computational
device (orthonormal bases), through orthogonal reframings that preserve spectra and subspace
determinants, through the averaging pushforward \(y\mapsto S^\dagger y\) of
\SG{cor:ig-full-rank-restriction-pushforward}, whose equality condition uses the Euclidean complement
and is chart-level under \(\GL(K)\), and through the finite part of the regularized identification cost,
whose divergent coefficient \((n-r)/2\) is frame free.

\section{Open ends}
\label{comp:sg1-open}

The supplement records three open problems here: the law of a retained mass coordinate under a scale
map, settled by an explicit induced law with a proof that its orbit stays in the natural domain or
reaches the boundary (\SG{open:exp-mass-sector}); whether the affine information projection commutes with
gauge actions or coarse parameter maps, settled by an intertwining theorem or a counterexample
(\SG{open:exp-projection-commutation}); and whether the Bayesian-renormalization information scale, built
on the Fisher metric of generative parameters, transfers to \(\Lambda\) read as the Fisher metric of a
recognition mean, settled by making the coarse mean a generative parameter with nondegenerate Fisher
metric and transferring the construction, or by a family where that metric is singular
(\SG{open:ig-model-space-fisher-transfer}). G5 declines any contraction claim for varying precision
(\SG{open:grg-galerkin-monotonicity}), and G4 is a Gaussian instance of
\MT{open:restrict-intrinsic-criterion} and \MT{open:cg-natural-selector}. Four gaps are visible but not recorded as open problems: law-level
membership for precisions whose graph has cycles; the relation of the family to a covariant Laplacian
with nontrivial holonomy; the combined mean-tie and factorization cost as a selector; and complex
Gaussians and complex or unitary frame groups beyond a real representation. These describe what the
documents do not contain; they are not proposals.

\section{Verdict}
\label{comp:sg1-verdict}

\begin{verdictbox}{\DELIVERED}
\textbf{What the part says it does.} Explicit Gaussian forms of the theory's objects: exponential-family
boundary and projection results; the information form, Schur lemma, and frame-covariant kernels of the
directed law; an exact account of when the belief precision is a matrix-weighted Laplacian plus self
terms, with that shape adopted as a hypothesis rather than derived; restriction costs, with the
statement that neither the factorization gap nor the mean cost selects a structure; and a separation of Fisher and spectral quantities into invariants and chart or
frame artifacts.

\textbf{What the statements give.} All of it, at arbitrary finite \(N\) and \(K\), with proofs. Much is
textbook in substance and presented as such. The program-specific content, the membership analysis of G3
and the invariance typing of G5, is elementary linear algebra done carefully; for \(K=1\) its sign
conditions are the Mattis-gauge and frustration conditions of spin glasses.

\textbf{The limits a reader must carry.} The interaction family is the pure-gauge sector, and the chart
in which it is imposed, not the law, decides what agreement means downstream. Law-level membership is settled on forests and pairs
and only bounded by a necessary sign condition on cycles, where a 512-agent network lives. The pencil
spectrum is data of a law with its link-fixed split, so the remark offering it as a frame-independent
membership criterion claims more than the propositions give. The combined recognition cost is not
examined as a selector. Absorption and the rank-only label need \(\GL^+(K)\) frames; under \(\GL(K,\R)\)
they hold and absorption is easier, compact groups keep more invariants and lose absorption, and complex
Gaussians are not treated. None of this is a defect in a proof, and the part states most of it itself.
\end{verdictbox}
